\documentclass[5pt]{article}
\usepackage{multicol,multirow}
\usepackage{graphicx} % Required for inserting images
\usepackage[margin=0.45cm]{geometry}
\usepackage{xcolor}
\usepackage{array}

\definecolor{LightGray}{gray}{0.9}

\newcolumntype{L}{>{\footnotesize}l}
\newcolumntype{C}{>{\ttfamily\footnotesize}l}

\usepackage{minted}

\begin{document}
\begin{center}
     \Large{\textbf{Julia Cheat Sheet}}\\
     \small{Book: Practical Julia: A hands-on introduction for scientific minds }\hfill\small{\textcopyright Maximilien Notz \the\year{}}
     \noindent\rule{20cm}{0.4pt}
\end{center}
\begin{multicols}{2}
\setcounter{secnumdepth}{0}


\section{Types and Variables}
Julia is a dynamically typed language, but it is possible to specify the type of a variable with the `::` operator.

\subsection{Lowlevel types}
\begin{tabular}{C L C L}
Int64 & 64-bit signed integer. & UInt64 & 64-bit unsigned integer.\\
Int32 & 32-bit signed integer. & UInt32 & 32-bit unsigned integer.\\
Int16 & 16-bit signed integer. & UInt16 & 16-bit unsigned integer.\\
Int8  & 8-bit signed integer. & UInt8  & 8-bit unsigned integer.\\
Float64 & 64-bit floating point. & Float32 & 32-bit floating point.\\
Float16 & 16-bit floating point. & Float8 & 8-bit floating point.\\
bool & Boolean type. & Char & Character type.\\
\end{tabular}

\subsection{Mathematical types}
The following types are abstract types, which means that they cannot be instantiated directly, but they can be used to define other types.
\begin{tabular}{C L}
    isequal(x,y) & Check if 2 different types variables are equal.\\
    abs(x) & Get the absolute value of \texttt{x}.\\
    abs2(x) & Get the absolute value of \texttt{x} squared.\\
    inf & Positive infinity constant.\\
\end{tabular}

\subsubsection{Reals}
\begin{tabular}{C L}
    Real\{T\} & Real number with type \texttt{T}.\\
\end{tabular}

\subsubsection{Complexes}
\begin{tabular}{C L}
    Complex\{T\} & Complex number with type \texttt{T} real and imaginary parts.\\
    im & Imaginary unit, e.g., \texttt{1.0im}.\\
    imag(z) & Get the imaginary part of the complex number \texttt{z}.\\
    real(z) & Get the real part of the complex number \texttt{z}.\\
    conj(z) & Get the complex conjugate of the complex number \texttt{z}.\\
    angle(z) & Get the phaseangle of the complex number \texttt{z}.\\
\end{tabular}

\subsubsection{Rationals}
\begin{tabular}{C L}
    Rational\{T\} & Rational number with type \texttt{T} numerator and denominator.\\    
    numerator(r) & Get the numerator of the rational number \texttt{r}.\\
    denominator(r) & Get the denominator of the rational number \texttt{r}.\\
\end{tabular}


\end{multicols}
\end{document}
